% Vietnamese translation for OI Public Library
% Source-File: IOI中国国家候选队论文/2009/贾志豪/组合游戏略述——浅谈SG游戏的若干拓展及变形.ppt
\documentclass[11pt]{article}
\usepackage{oipl-vn}

\newcommand{\mex}{\operatorname{mex}}

\title{Lược khảo trò chơi tổ hợp: Một số mở rộng và biến thể của trò chơi SG\\{\large Ghi chú theo slide}}
\author{Jia Zhihao\\Trường Trung học số 2 Thạch Gia Trang}
\date{2009}

\begin{document}
\maketitle

\origtitle{Extensions and variants of SG games}
\sourcepath{IOI-China-Candidate-Team-Papers/2009/Jia-Zhihao/sg-game-variants-slides.ppt}

\section{Nền tảng}

Slide bắt đầu từ trò chơi công bằng, hữu hạn, không có yếu tố may rủi. Mỗi
trạng thái là một đỉnh trong đồ thị có hướng; nước đi là cạnh. Hàm SG của một
đỉnh là \(\mex\) của các giá trị SG ở các đỉnh kề.

\section{Các biến thể luật}

\begin{itemize}
  \item Anti-SG: người đi nước cuối thua, nên điều kiện thắng thua của tổng
  trò chơi thay đổi so với chuẩn SG.
  \item Multi-SG: một đống có thể tách thành nhiều đống, làm một nước đi sinh
  ra tổng của nhiều trò con.
  \item Every-SG: mọi quân có nước đi đều phải đi, nên thời gian kết thúc của
  từng thành phần trở thành thông tin quan trọng.
\end{itemize}

\section{Mô hình thường gặp}

Trò lật đồng xu, xóa cạnh trên cây, \textit{Christmas Game} và xóa cạnh trên
đồ thị vô hướng cho thấy mỗi biến thể cần được kiểm tra lại từ luật chơi, thay
vì áp dụng máy móc công thức xor.

\section{Ghi chú}

Bản slide là bảng phân loại nhanh. Khi gặp bài mới, trước hết xác định nó có
đúng là trò chơi SG chuẩn hay không; nếu có luật khác về nước cuối, tách đống
hoặc di chuyển đồng thời, cần dùng biến thể tương ứng.

\end{document}
